% 第3章3.5.2匹配的单位容量流网络：源 s 到左侧 a,b,c，允许边 a-x,a-y,b-y,b-z,c-y 改为左到右的弧，右侧 x,y,z 到汇 t，容量均为1。
% 匹配 {a-x,b-y} 对应流值2的整数流（蓝）；残量路径 s->c->y->b->z->t（绿）中 y->b 是撤回匹配边 b-y 的反向残量弧。
\begin{tikzpicture}[>=Stealth,line width=.8pt,font=\fontsize{8}{10}\selectfont,
  v/.style={circle,draw=FigureStroke,fill=FigureBlueFill,line width=.8pt,minimum size=6.5mm,inner sep=0pt,text=FigureInk},
  e/.style={->,draw=FigureMuted,line width=.8pt,shorten >=.5pt},
  f/.style={->,draw=FigureBlue,line width=1.6pt,shorten >=.5pt},
  p/.style={->,draw=FigureGreen,line width=1.4pt,dashed,shorten >=.5pt},
  note/.style={text=FigureMuted,font=\FigureNoteFont\fontsize{7.5}{9}\selectfont}]
  \node[v] (s) at (0,0) {$s$};
  \node[v] (a) at (1.9,1.2) {$a$};
  \node[v] (b) at (1.9,0) {$b$};
  \node[v] (c) at (1.9,-1.2) {$c$};
  \node[v] (x) at (4.3,1.2) {$x$};
  \node[v] (y) at (4.3,0) {$y$};
  \node[v] (z) at (4.3,-1.2) {$z$};
  \node[v] (t) at (6.2,0) {$t$};
  \draw[f] (s)--(a); \draw[f] (s)--(b); \draw[e] (s)--(c);
  \draw[f] (a)--(x); \draw[e] (a)--(y); \draw[f] (b)--(y); \draw[e] (b)--(z); \draw[e] (c)--(y);
  \draw[f] (x)--(t); \draw[f] (y)--(t); \draw[e] (z)--(t);
  % 残量路径
  \draw[p] (s) to[bend right=22] (c);
  \draw[p] (c) to[bend right=16] (y);
  \draw[p] (y) to[bend right=22] (b);
  \draw[p] (b) to[bend right=16] (z);
  \draw[p] (z) to[bend right=22] (t);
  \node[note] at (3.1,1.85) {反向残量弧$y\to b$表示撤回匹配边$b-y$};
  \node[note] at (3.1,-2.05) {所有弧容量为$1$；蓝色粗线：当前流（匹配$\{a-x,b-y\}$）；绿色虚线：残量增广路径};
\end{tikzpicture}
